%% ma: L12-B17-0006
%% lop: 12
%% chuong: 5
%% bai: 17
%% dang: 12.17.6
%% loai: TN4
%% muc: TH
%% dap_an: "A"
%% nguon: "(NBV)-ĐỀ SỐ 7-MỖI NGÀY 1 ĐỀ THI 2025.pdf (D07, MỖI NGÀY 1 ĐỀ - Đề số 7), Phần I Câu 11"
%% kien_thuc: "Bài 17 (phương trình mặt cầu biết tâm và một điểm thuộc mặt cầu)"
%% trang_thai: cho-duyet
%% ly_do: "Dạng toán mới đề xuất 12.17.6 (Phương trình mặt cầu biết tâm và một điểm thuộc mặt cầu); chờ giáo viên duyệt dạng."
%% ghi_chu: "Đáp án gốc A đúng."
\begin{ex}
Trong không gian $Oxyz$, cho hai điểm $I(1;-2;1)$ và $A(1;2;3)$. Phương trình mặt cầu có tâm $I$ và đi qua $A$ là
\choice{\True $(x-1)^2+(y+2)^2+(z-1)^2=20$}{$(x+1)^2+(y-2)^2+(z+1)^2=5$}{$(x+1)^2+(y-2)^2+(z+1)^2=20$}{$(x-1)^2+(y+2)^2+(z-1)^2=5$}
\loigiai{Bán kính $R=IA=\sqrt{0^2+4^2+2^2}=\sqrt{20}$, nên $R^2=20$. Phương trình mặt cầu: $(x-1)^2+(y+2)^2+(z-1)^2=20$.}
\end{ex}
